\section{楕円gluingの導来比較と全heartの一回のtilt}
\label{sec:resonance-global-heart}
本節では，古典的avatarがgenuineである楕円gluingに範囲を限定し，
自然な導来切断の完全忠実性を使えることを確定する．
導来比較自体は現行の別稿
\texttt{papers/dp1\_anticanonical\_models.tex}の主定理であり，本節の証明はその論理と入力を改めて確認したものである．
ここではその証明を，orbit生成，traceの保存，飽和性の三段階に分けて
点検し，必要な一次資料の適用条件を明示する．
そこから得る全heartの振幅，torsion pair，および混合$A_2$の
アフィン表現の結論を本稿の結果に組み込む．
以下の結果を任意の非可換del Pezzo曲面，特異楕円境界，
非局所自由gluing，あるいはweak avatarへ無条件に適用しない．

\subsection{比較定理の入力と証明}
$S$を滑らかな可換del Pezzo曲面，$E\in|-K_S|$を滑らかな曲線，
$q\in\Pic^0(E)$とする．$\mathcal B=\mathcal O_S^\perp\subset\Perf(S)$に
exceptional object $e$を
\[
\RHom(e,M)=0,\qquad
\RHom(M,e)=\RHom_E(Li^*M,q)\quad(M\in\mathcal B)
\]
でglueして得るdg圏を$\mathcal D_q$とする．Rains heartを
$\mathcal A_q$とし，$\mathcal D_q=D^b(\mathcal A_q)$である．
境界関手$F_q$の規約は
\[
F_q(e)=\mathcal O_E,\qquad F_q(M)=Li^*M\otimes q^{-1}
\quad(M\in\mathcal B)
\]
である．このgluing，境界triangle，$\theta=\mathbb S[-2]$のheart保存は
Rains\cite[\S4, (4.4), (4.7)]{RainsFiltered}および
\cite[Theorem 4.30, \S5]{RainsBirational}の入力である．
$\tau=\theta^{-1}$，$R=\bigoplus_{n\ge0}\Hom(e,\tau^ne)$，
$\mathcal O_R(a)=\pi R(a)$とする．

\begin{theorem}[genuine楕円gluingにおける指定導来比較]
\label{thm:resonance-global-comparison}
上の設定で，実際のRains heartのinjective resolutionから作る
自然なdg関手
\[
\Phi=\pi\mathbf R\Gamma_*:
\mathcal D_q\longrightarrow D^b(\qgr R),\qquad
(\mathbf R\Gamma_*X)_n=\RHom(e,\tau^nX)
\]
は同値であり，$\Phi(e)=\mathcal O_R$，$\Phi\tau=(1)\Phi$を満たす．
$q$または境界平行移動の位数に制限はない．
また$\qgr R$のhomological dimensionは2である．
これは異なる$q$間の圏が同値であるという主張ではない．
\end{theorem}
\begin{proof}
\emph{環とorbit．}
Rains\cite[Proposition 5.1]{RainsFiltered}より，$R$は中心正則元
$g\in R_1$を持ち，$R/gR=B(E,L,\sigma)$，$\deg L=d=(-K_S)^2>0$．
$R$は$A=R/(g-1)$のRees環である．
$\operatorname{gr}_{(g)}R=B[g]$が整域で，固定内部次数では$g$進filtrationが有限なので，
$R$も整域である．同じinitial termの持上げから左右noetherian性を得る．
$\sigma$が無限位数ならRogalski--Sierra--Stafford
\cite[Proposition 2.4]{RSSResonance}が，$\deg L\ge1$のままで
左右の$\chi$と有限cohomological dimensionを与える．
有限位数なら$B$は$k(E)[t,t^{-1};\sigma]$の部分環なのでPIであり，
Yekutieli--Zhang\cite[Theorems 5.13, 5.1(1), 4.5]{YZResonance}を
この順に使って同じ有限性を得る．ここで必要なのは
$\mathcal O_R(a)$からのExtの有界性であって，標的の大域次元ではない．

sourceの境界triangle，$e$のexceptionality，楕円曲線上の
正次数線束の$H^1$消滅およびSerre双対性から
\begin{equation}\label{eq:comparison-orbit}
\Ext^j(\tau^ae,\tau^be)=
\begin{cases}
 R_{b-a},&j=0,\ b\ge a,\\
 R_{a-b-1}^*,&j=2,\ b<a,\\
 0,&\text{それ以外}
\end{cases}
\end{equation}
を得る．同じ式が$\mathcal O_R(a),\mathcal O_R(b)$にも成り立ち，
そのHom同定は自然な飽和写像である．標的の入力は
Rains\cite[Corollary 3.8とその証明]{RainsFiltered}である．
同系の次数2のpairingは境界$H^1(E,\mathcal O_E)$へのtraceで定まる．
この段階で標的全体のsmoothnessやSerre関手を先取りしない．

\emph{sourceのthick生成．}
$X\in\operatorname{Ind}(\mathcal D_q)$が全$\tau^ne$に右直交するとする．
$\theta X$も右直交なので，境界triangleと随伴性により
\[
\RHom_E(\Psi^n\mathcal O_E,F_qX)=0\quad(n\in\mathbb Z),
\qquad F_q\tau=\Psi F_q
\]
となる．正次数のtwisted line bundlesは$D_{\rm qc}(E)$を生成するため
$F_qX=0$．また$\RHom(e,X)=0$から$X\in\operatorname{Ind}(\mathcal B)$で，
対応する$G\in D_{\rm qc}(S)$は$R\Gamma(G)=0$，$Li^*G=0$を満たす．
$U=S\setminus E$は$E$がampleなのでaffineである．$j:U\hookrightarrow S$に対し
\[
G\simeq\operatorname*{hocolim}_{n}G(nE)\simeq Rj_*j^*G
\]
であり，affine上の大域切断の保守性から$G=0$となる．
ゆえにInd圏でorbitがcompact生成し，そのcompact objectsは
$\mathcal D_q=\operatorname{thick}\{\tau^ne:n\in\mathbb Z\}$である．
compact objectsの中だけで直交圏を計算して生成性を推論してはいない．

\emph{自然関手とtrace．}
exactな$\tau$と実際のheartの射から次数付き$R$作用を持つ
$\mathbf R\Gamma_*$を構成する．式\eqref{eq:comparison-orbit}より
$\Phi(\tau^ae)=\mathcal O_R(a)$であり，高次の有限内部次数は$\pi$で消える．
thick生成により全$X$の像が$D^b(\qgr R)$に入る．
境界対象では随伴性から$R\Gamma(E,\Psi^nV)$となり，Serre消滅により
$\Phi(i_*V)$は$\qgr B\simeq\coh E$の通常の像である．
この同定は短完全列を保つ．
非零$u\in H^1(E,\mathcal O_E)$を$\mathcal O_E$の自己拡大で表し，
$e\to i_*\mathcal O_E$で引き戻して境界triangleのconnecting mapを取る．
$e$と$\mathcal O_R$のexceptionalityにより，両側で
\[
H^1(E,\mathcal O_E)\simeq\Ext^2(e,\theta e),\qquad
H^1(E,\mathcal O_E)\simeq\Ext^2(\mathcal O_R,\mathcal O_R(-1))
\]
となり，$\Phi$はこの$u$を同じ$u$へ送る．
従ってtraceの写像は非零である．全shiftとYoneda積を使うと，
全orbit間で次数2の写像は次数0の同型の双対となる．
式\eqref{eq:comparison-orbit}とthick生成を両変数に適用して
$\Phi$の完全忠実性を得る．Extの次元一致だけを使う議論ではない．

\emph{本質的全射性．}
$\mathcal B$は$\Perf(S)$のadmissible部分圏なのでsmooth properである．
gluing moduleは各対象上で有限次元の全cohomologyを持ち，
Orlovのperfectness判定によりperfectである．従ってsourceはsmooth properであり，
Orlov\cite[Theorems 3.18, 3.25, Corollary 3.19]{OrlovResonance}から飽和的である．
任意の$Y\in D^b(\qgr R)$に対し
$X\mapsto\RHom(\Phi X,Y)$は，環の有限性とsourceのthick生成から
pseudo-perfect moduleとなり，sourceの対象で表現される．
従って$\Phi$の右随伴が存在する．counitのconeは全
$\mathcal O_R(a)$に右直交する．標的のample pair，$\chi$，有限cohomological
 dimensionから，その有限個のcohomology sheavesの十分高いtwistに
Serre消滅を同時に適用すると，非零cohomologyがある場合の
非零大域切断と矛盾する．よってconeは零であり全射性を得る．

最後にsourceのSerre関手を同値で運ぶと標的全体で$(-1)[2]$となる．
$(-1)$がheartを保つためExtは次数2を超えて消え，
$\Ext^2(\mathcal O_R,\mathcal O_R(-1))=k$からhomological dimensionは2である．
\end{proof}

\subsection{全heartでの振幅とtorsion pair}
\begin{theorem}[全Rains heartの一回のHRS tilt]
\label{thm:resonance-global-tilt}
定理\ref{thm:resonance-global-comparison}の設定で，
$H=-K_X$がnefであるRains heartを取る．
全$M\in\mathcal A_q$について$\Phi(M)$のcohomological amplitudeは$[0,1]$である．
\[
\mathcal T=\{M:H^1\Phi(M)=0\},\qquad
\mathcal F=\{M:H^0\Phi(M)=0\}
\]
は$\mathcal A_q$上のtorsion pairで，
\[
\boxed{\Phi^{-1}(\qgr R)=\langle\mathcal F[1],\mathcal T\rangle.}
\]
$q$が非torsionなら境界外Serre部分圏$\mathcal Z$へ制限でき，そこで
\[
\mathcal T\cap\mathcal Z=\ker\Ext^1(e,-),\qquad
\mathcal F\cap\mathcal Z=\ker\Hom(e,-)
\]
となる．このtorsion pairは全$\mathcal Z$に対するものである．
\end{theorem}
\begin{proof}
Rains heartはline bundlesで生成される
\cite[Theorem 9.10]{RainsBirational}ので，
$M$への全射$\bigoplus_{j=1}^sL_j\twoheadrightarrow M$を取れる．
Serre双対性により
\[
\Ext^2(e,\tau^nM)=\Hom(M,\theta^{n+1}e)^*.
\]
右辺の双対は$\bigoplus_j\Hom(L_j,\theta^{n+1}e)$へ単射する．
非零射があれば，rank-one torsion-free sheaf間の射に対する
Rains\cite[Corollary 10.12]{RainsBirational}から
$(n+1)K_X-c_1(L_j)$は有効でなければならない．
しかしその$H$との交点は$-(n+1)d-Hc_1(L_j)<0$ $(n\gg0)$であり，
nef性に反する．よって$\Ext^2(e,\tau^nM)=0$ $(n\gg0)$．
$\mathbf R\Gamma_*M$の次数2のcohomologyは有限内部次数にしかなく，
各成分が有限次元なので$\pi$で零になる．その他の次数はsourceのhdim2と
heart性から消え，振幅$[0,1]$を得る．

$\mathcal B'=\Phi^{-1}(\qgr R)$とすると，有界t構造の直交性から
$\mathcal B'\subset D^{[-1,0]}(\mathcal A_q)$である．
中間t構造のHRS対応により，$\mathcal T=\mathcal A_q\cap\mathcal B'$，
$\mathcal F=\mathcal A_q\cap\mathcal B'[-1]$がtorsion pairで，
$\mathcal B'=\langle\mathcal F[1],\mathcal T\rangle$となる．
具体的には，$M\in\mathcal A_q$の$\mathcal B'$-truncation triangleを
$\mathcal A_q$で読むと$0\to T\to M\to F\to0$を与え，
負次数のHom消滅から$\Hom(\mathcal T,\mathcal F)=0$を得る．
$\mathcal Z$はSerre部分圏なので$T,F$もそこに入り，
その上では$g$周期性が$H^j\Phi(M)=0\iff\Ext^j(e,M)=0$を与える．
\end{proof}
\begin{remark}
全heartに対する定義は$H^j\Phi$によるものである．
境界外でない任意の対象に対し，有限生成性を点検せず生の次数付きExtの
各成分の消滅へ置き換えない．また本定理はVdB/AS heartの指定対象との
同定や，相対モジュライへの降下を証明するものではない．
\end{remark}

\subsection{境界外の導来Extと二つのcharacter}
\begin{lemma}[中心局所化による導来Extの保存]
\label{lem:resonance-affine-ext}
$A=(R[g^{-1}])_0=R/(g-1)$の有限次元加群$U,V$を
境界外の$\qgr R$の対象とみなすと，全$j\ge0$で
\[
\Ext^j_A(U,V)\simeq\Ext^j_{\qgr R}(U,V)
\]
となる．
\end{lemma}
\begin{proof}
$R[g^{-1}]$は$g$を可逆元に持つ強次数付き環であり，
その次数付き加群圏は$\operatorname{Mod}A$と同値である．
中心局所化はexactで，その右随伴はrestrictionとtorsion quotientでありexact．
局所化加群は内部次数torsionに導来的にも直交するため，この随伴は
$\operatorname{QGr}R$でも完全忠実である．従って導来随伴も完全忠実となる．
有限次元$U$に対応する$g$可逆次数付き加群の非負次数部分は
$k[g]$上有限自由であり，有限生成$R$格子を与える．ゆえにその像は
$\qgr R$に属する．locally noetherian圏でのcoherent Extとの一致を使えば結論を得る．
\end{proof}

\begin{theorem}[混合$A_2$のcharacter分離と非分裂三角表現]
\label{thm:resonance-affine-triangular}
上のgenuine楕円gluingで，指数$+2,-1$の実在混合根対を取る．
$V=\Phi(S_+)$，$U=\Phi(S_-)[1]$，$W=\Phi(P)$とする．
これらはそれぞれ次元$2,1,1$の$A$加群で，
\[
0\longrightarrow W\longrightarrow V\longrightarrow U\longrightarrow0
\]
は非分裂であり，$W\not\simeq U$である．従って$A$の$V$への作用像は
基底の選択を除き全下三角$2\times2$行列代数である．
また$\Hom_A(V,U)=k$である．
\end{theorem}
\begin{proof}
定理\ref{thm:a2-delta-rank}から表示した短完全列を得る．
既存のHom表とambient Euler形式から$\Ext_X^1(P,S_-)=0$である．
導来比較と補題\ref{lem:resonance-affine-ext}により
$\Hom_A(W,U)=\Ext_X^1(P,S_-)=0$，従って$W\not\simeq U$．
非零包含$S_-\hookrightarrow P$は比較同値によって非零射
$U[-1]\to W$へ移る．そのconeのtriangleは上の加群拡大であるから
拡大類は非零である．
二つの異なるcharacterを持つ二次元非分裂表現の像には，
異なる対角成分を持つ元とその多項式による二つの冪等元がある．
それらに適合する基底では非分裂性が非零の非対角成分を与えるので，
作用像は次元3の全三角行列代数となる．最後のHomはこの表示から従う．
\end{proof}

\begin{theorem}[二つのcharacterが定める四次元商]
\label{thm:resonance-preprojective-quotient}
上の$W,U$が生成する$\operatorname{mod}_{\rm fd}A$のSerre部分圏を
$\mathcal D$とする．$\Ext_A^1(W,U)=k$の非零元による
$0\to U\to V'\to W\to0$を取ると，
\[
\mathcal D\simeq\operatorname{mod}_{\rm fd}\Pi(A_2),\qquad
I=\operatorname{Ann}_A(V\oplus V'),\qquad
\boxed{A/I\simeq\Pi(A_2).}
\]
$\mathcal D$の直既約は$W,U,V,V'$の四つである．
この商の次元は4で，そのJacobson radicalの平方は零である．
\end{theorem}
\begin{proof}
導来比較と既知のHom/Euler表から，$W,U$はともに2-sphericalであり，
両向きのcross HomとExt$^2$は零，cross Ext$^1$は一次元である．
例えば$\Ext_A^1(W,U)=\Ext_X^2(P,S_-)=\Hom_X(S_-,P)^*=k$であり，
逆方向は$\Hom_X(S_-,P)=k$から得る．
境界外では自然にSerre関手が$[2]$となるので，
非分裂の両拡大$V,V'$のHomとExt$^1$は先の二根の計算と同じである．
従って$V\oplus V'$は$\mathcal D$のprojective generatorで，
End代数は$\Pi(A_2)$，直既約は表示した四つとなる．

商環の主張には単なるMorita同値以上の確認が必要である．
$W$と$U$のcharacterをそれぞれ$\chi_W,\chi_U$とすると，
$V\oplus V'$への像は，二つの対角成分を共有し，二つの独立な
非対角成分を持つ四次元代数$C_0\simeq\Pi(A_2)$の部分代数$C$に入る．
$\chi_W\ne\chi_U$より$C$の対角像は$k\times k$全体．
$J=\operatorname{rad}C_0$は$J^2=0$なので，対角冪等元は$C$へ持ち上がる．
$1+J$による共役でそれらを$C_0$の標準冪等元へ移す．
両拡大の非分裂性により両非対角cornerが非零である．
各cornerは一次元なので$C=C_0$となる．
従って核$I$による商が$\Pi(A_2)$である．
\end{proof}
\begin{corollary}[二根部分系におけるWeyl反射]
\label{cor:resonance-weyl-reflection}
導来比較で$K$群を同定すると，二つの一次元加群の根類は
\[
c_1(W)=\alpha+\beta=s_\beta(\alpha),\qquad
c_1(U)=-\beta=s_\beta(\beta),
\quad s_\beta(v)=v+(v\beta)\beta.
\]
従ってこの二根部分系では，指数$(2,-1)$から新しい単純根の指数$(1,1)$への
変化はWeyl反射で記述される．これは全$E_6$のheartを一つの反射で記述する主張ではない．
\end{corollary}
\begin{proof}
$[P]=[S_+]+[S_-]$，$[S_-[1]]=-[S_-]$から根類の式が従う．
$\alpha\beta=1$，$\beta^2=-2$を反射の定義へ代入し，
$\chi(e,P)=1$，$\chi(e,S_-[1])=1$を使う．
\end{proof}

\begin{remark}[決定した表示と未計算の係数]
この結果は$A$の有限次元商と表現の同型型を決定するが，
元の楕円データで指定された$A$の各生成元に対する
$\chi_W(a),\chi_U(a),\eta(a)$の数値式はまだ与えていない．
また$A\twoheadrightarrow\Pi(A_2)$はhomological epimorphismではない．
商の単純加群には二周期的な高次Extがあるが，$A$上では
補題\ref{lem:resonance-affine-ext}とhdim2からExt$^{>2}$は零だからである．
この区別は，有限長heartの同定からambientの導来圏を同定する際に不可欠である．
\end{remark}
